\documentclass{article}
\usepackage{amsmath,amssymb,amsthm}
\usepackage{algorithm}
\usepackage{algorithmic}
\usepackage{enumitem}
\usepackage{hyperref}
\usepackage{geometry}
\geometry{margin=1in}
\newtheorem{theorem}{Theorem}
\newtheorem{lemma}{Lemma}
\newtheorem{remark}{Remark}
\newcommand{\prox}{\operatorname{prox}}
\newcommand{\E}{\mathbb{E}}
\newcommand{\R}{\mathbb{R}}

\begin{document}

\section*{Algorithm Name}
\textbf{Proximal Gradient Method with Gradient Momentum (PGM-GM)}  

The method solves the composite convex problem  
\[
\min_{x\in\R^{d}}\; \Phi(x):= f(x)+g(x),
\]
where $f$ is smooth and convex and $g$ is proper, closed and convex (possibly nonsmooth).

---------------------------------------------------------------------

\section*{Mathematical Formulation}
Given a stepsize $\eta>0$ and a momentum parameter $\beta\in[0,1)$, the iterates are defined by  
\begin{align}
v_k &:= \nabla f(x_k)+\beta\bigl(\nabla f(x_k)-\nabla f(x_{k-1})\bigr), \qquad k\ge 1, \label{eq:momentum}\\[4pt]
x_{k+1} &:= \prox_{\eta g}\!\bigl(x_k-\eta v_k\bigr), \qquad k\ge 0, \label{eq:update}
\end{align}
with an initial point $x_0\in\R^{d}$ and a dummy gradient $\nabla f(x_{-1})\equiv0$ (so that $v_0=\nabla f(x_0)$).

The proximal operator is
\[
\prox_{\eta g}(z):=\arg\min_{u\in\R^{d}}\Bigl\{ g(u)+\frac{1}{2\eta}\|u-z\|^{2}\Bigr\}.
\]

---------------------------------------------------------------------

\section*{Pseudocode}
\begin{algorithm}[H]
\caption{PGM-GM}
\begin{algorithmic}[1]
\STATE \textbf{Input:} stepsize $\eta>0$, momentum $\beta\in[0,1)$, max iterations $T$
\STATE Choose $x_0\in\R^{d}$, set $\nabla f(x_{-1})\gets0$
\FOR{$k=0,\dots,T-1$}
    \STATE Compute gradient $g_k\gets \nabla f(x_k)$
    \IF{$k=0$}
        \STATE $v_k\gets g_k$ \COMMENT{no momentum at first step}
    \ELSE
        \STATE $v_k\gets g_k+\beta\bigl(g_k-g_{k-1}\bigr)$ \label{alg:momentum}
    \ENDIF
    \STATE $x_{k+1}\gets \prox_{\eta g}\!\bigl(x_k-\eta v_k\bigr)$ \label{alg:update}
\ENDFOR
\STATE \textbf{Output:} $x_T$
\end{algorithmic}
\end{algorithm}

---------------------------------------------------------------------

\section*{Dry Run Example}
Consider the scalar problem  
\[
f(w)=(w-3)^2,\qquad g\equiv0,\qquad \eta=0.1,\qquad \beta=0.5,\qquad w_0=0.
\]
Because $g\equiv0$, $\prox_{\eta g}= \operatorname{id}$.

\begin{center}
\begin{tabular}{c|c|c|c|c}
$k$ & $w_k$ & $\nabla f(w_k)=2(w_k-3)$ & $v_k$ (momentum) & $w_{k+1}=w_k-\eta v_k$ \\ \hline
0 & $0$ & $-6$ & $v_0=-6$ & $w_1=0-0.1(-6)=0.6$ \\[2pt]
1 & $0.6$ & $2(0.6-3)=-4.8$ & $v_1=-4.8+0.5(-4.8+6)= -4.8+0.5(1.2)=-4.2$ &
$w_2=0.6-0.1(-4.2)=1.02$ \\[2pt]
2 & $1.02$ & $2(1.02-3)=-3.96$ & $v_2=-3.96+0.5(-3.96+4.8)= -3.96+0.5(0.84)= -3.54$ &
$w_3=1.02-0.1(-3.54)=1.374$ 
\end{tabular}
\end{center}
Objective values: $f(w_0)=9$, $f(w_1)=5.76$, $f(w_2)=3.9204$, $f(w_3)=2.777$ – the method reduces the objective.

---------------------------------------------------------------------

\section*{Assumptions}
\begin{enumerate}[label=(A\arabic*)]
\item \label{A1} \textbf{Convexity of $f$ and $g$:}  
$f:\R^{d}\to\R$ and $g:\R^{d}\to\R\cup\{+\infty\}$ are convex.
\item \label{A2} \textbf{Smoothness of $f$:}  
$f$ is $L$‑smooth, i.e.
\[
\|\nabla f(x)-\nabla f(y)\|\le L\|x-y\|,\qquad\forall x,y\in\R^{d}.
\]
Equivalently,
\[
f(y)\le f(x)+\langle\nabla f(x),y-x\rangle+\frac{L}{2}\|y-x\|^{2}. \tag{1}
\]
\item \label{A3} \textbf{Existence of a minimizer:}  
There exists $x^{\star}\in\arg\min_{x}\Phi(x)$ and the optimal value $\Phi^{\star}:=\Phi(x^{\star})$ is finite.
\item \label{A4} \textbf{Stepsize bound:}  
\[
0<\eta\le\frac{1}{L(1+\beta)}.
\]
\end{enumerate}

---------------------------------------------------------------------

\section*{Main Convergence Theorem}
\begin{theorem}\label{thm:main}
Under Assumptions \ref{A1}–\ref{A4}, the iterates $\{x_k\}$ generated by \eqref{eq:momentum}–\eqref{eq:update} satisfy for all $K\ge 1$
\[
\Phi\!\bigl(\bar{x}_{K}\bigr)-\Phi^{\star}
\;\le\;
\frac{\|x_{0}-x^{\star}\|^{2}}{2\eta K},
\qquad\text{where }\;
\bar{x}_{K}:=\frac{1}{K}\sum_{k=0}^{K-1}x_{k}.
\]
Consequently $\Phi(\bar{x}_{K})\to\Phi^{\star}$ at the rate $O(1/K)$.
\end{theorem}

---------------------------------------------------------------------

\section*{Theoretical Properties}
\begin{itemize}
\item \textbf{Stability:} The bound on $\eta$ guarantees that the ``effective'' gradient $v_k$ is $L(1+\beta)$‑Lipschitz, preventing divergence even when $\beta$ is close to $1$.
\item \textbf{Hyperparameter Sensitivity:} Larger $\beta$ accelerates early progress but forces a smaller admissible stepsize by \ref{A4}. The $O(1/K)$ rate is independent of $\beta$ as long as the stepsize condition holds.
\item \textbf{Comparison to Baselines:}
    \begin{enumerate}
    \item With $\beta=0$, the method reduces to the classical proximal gradient method, recovering the standard $O(1/K)$ rate.
    \item Compared to Nesterov’s accelerated proximal gradient, PGM‑GM is simpler (only a momentum on the gradient) and enjoys the same worst‑case rate for the convex case, while often exhibiting better practical stability.
    \end{enumerate}
\end{itemize}

---------------------------------------------------------------------

\section*{Discussion}
The proof follows the classic proximal‑gradient analysis but incorporates the extra momentum term. The key is to view $v_k$ as a perturbed gradient whose perturbation can be bounded thanks to smoothness and the stepsize restriction. The resulting recursion is identical to that of the vanilla proximal method, leading to the same $O(1/K)$ guarantee.

---------------------------------------------------------------------

\section*{Preliminaries (Definitions and Lemmas)}
\begin{lemma}[Three‑point identity for proximal operator]\label{lem:prox}
For any $z\in\R^{d}$ and any $u\in\R^{d}$,
\[
\bigl\langle \frac{z-\prox_{\eta g}(z)}{\eta},\; \prox_{\eta g}(z)-u \bigr\rangle
\;\ge\;
g\!\bigl(\prox_{\eta g}(z)\bigr)-g(u).
\]
\end{lemma}
\begin{proof}
Standard; see e.g. \cite{parikh2014proximal}. \hfill $\square$
\end{proof}

\begin{lemma}[Descent Lemma for $L$‑smooth convex $f$]\label{lem:descent}
For any $x,y\in\R^{d}$,
\[
f(y)\le f(x)+\langle\nabla f(x),y-x\rangle+\frac{L}{2}\|y-x\|^{2}. 
\]
\end{lemma}
\begin{proof}
Directly from Assumption \ref{A2}. \hfill $\square$
\end{proof}

---------------------------------------------------------------------

\section*{Main Proof (Step‑by‑Step Derivation)}
We prove Theorem \ref{thm:main}. Throughout the proof we fix an arbitrary optimal point $x^{\star}\in\arg\min\Phi$.

\subsection*{Step 1 – Optimality relation from the proximal step}
From \eqref{eq:update} and Lemma \ref{lem:prox} with $z:=x_k-\eta v_k$ and $u:=x^{\star}$ we obtain
\[
\Bigl\langle v_k,\; x_{k+1}-x^{\star}\Bigr\rangle
\;\ge\;
g(x_{k+1})-g(x^{\star})+\frac{1}{2\eta}\Bigl(\|x^{\star}-x_k\|^{2}
-\|x^{\star}-x_{k+1}\|^{2}
-\|x_{k+1}-x_k\|^{2}\Bigr). \tag{2}
\]

\subsection*{Step 2 – Replace $v_k$ by the momentum expression}
Using the definition \eqref{eq:momentum},
\[
v_k = \nabla f(x_k)+\beta\bigl(\nabla f(x_k)-\nabla f(x_{k-1})\bigr).
\]
Insert this into the left‑hand side of (2):
\[
\langle\nabla f(x_k),x_{k+1}-x^{\star}\rangle
+\beta\langle\nabla f(x_k)-\nabla f(x_{k-1}),x_{k+1}-x^{\star}\rangle
\;\ge\;
\text{RHS of (2)}. \tag{3}
\]

\subsection*{Step 3 – Bounding the momentum cross term}
By Cauchy–Schwarz and the $L$‑smoothness (\ref{A2}),
\[
\|\nabla f(x_k)-\nabla f(x_{k-1})\|
\le L\|x_k-x_{k-1}\|.
\]
Hence,
\[
\beta\langle\nabla f(x_k)-\nabla f(x_{k-1}),x_{k+1}-x^{\star}\rangle
\le
\beta L\|x_k-x_{k-1}\|\;\|x_{k+1}-x^{\star}\|
\le
\frac{\beta L}{2}\bigl(\|x_k-x_{k-1}\|^{2}+\|x_{k+1}-x^{\star}\|^{2}\bigr). \tag{4}
\]

\subsection*{Step 4 – Apply the descent lemma to $f$}
From Lemma \ref{lem:descent} with $y:=x_{k+1}$,
\[
f(x_{k+1})\le f(x_k)+\langle\nabla f(x_k),x_{k+1}-x_k\rangle+\frac{L}{2}\|x_{k+1}-x_k\|^{2}. \tag{5}
\]

\subsection*{Step 5 – Combine (3)–(5)}
Rearrange (3) to isolate $\langle\nabla f(x_k),x_{k+1}-x^{\star}\rangle$:
\[
\langle\nabla f(x_k),x_{k+1}-x^{\star}\rangle
\ge
\text{RHS of (2)}-\beta\langle\nabla f(x_k)-\nabla f(x_{k-1}),x_{k+1}-x^{\star}\rangle.
\]
Insert the bound (4) and the expression of RHS of (2) to obtain
\begin{align}
\langle\nabla f(x_k),x_{k+1}-x^{\star}\rangle
&\ge
g(x_{k+1})-g(x^{\star}) \nonumber\\
&\quad+\frac{1}{2\eta}\bigl(\|x^{\star}-x_k\|^{2}
-\|x^{\star}-x_{k+1}\|^{2}
-\|x_{k+1}-x_k\|^{2}\bigr) \nonumber\\
&\quad-\frac{\beta L}{2}\bigl(\|x_k-x_{k-1}\|^{2}+\|x_{k+1}-x^{\star}\|^{2}\bigr). \tag{6}
\end{align}

\subsection*{Step 6 – Relate $\Phi(x_{k+1})-\Phi^{\star}$}
Add $f(x_k)-f(x^{\star})$ to both sides of (6) and use convexity of $f$:
\[
f(x_k)-f(x^{\star})\le\langle\nabla f(x_k),x_k-x^{\star}\rangle.
\]
Thus,
\begin{align}
\Phi(x_{k+1})-\Phi^{\star}
&=f(x_{k+1})-f(x^{\star})+g(x_{k+1})-g(x^{\star}) \nonumber\\
&\le \langle\nabla f(x_k),x_{k+1}-x^{\star}\rangle
+\frac{L}{2}\|x_{k+1}-x_k\|^{2}+g(x_{k+1})-g(x^{\star}) \nonumber\\
&\le \frac{1}{2\eta}\bigl(\|x^{\star}-x_k\|^{2}
-\|x^{\star}-x_{k+1}\|^{2}
-\|x_{k+1}-x_k\|^{2}\bigr) \nonumber\\
&\quad+\frac{L}{2}\|x_{k+1}-x_k\|^{2}
-\frac{\beta L}{2}\bigl(\|x_k-x_{k-1}\|^{2}+\|x_{k+1}-x^{\star}\|^{2}\bigr). \tag{7}
\end{align}

\subsection*{Step 7 – Use the stepsize bound}
Because $\eta\le\frac{1}{L(1+\beta)}$, we have
\[
\frac{L}{2}\le\frac{1}{2\eta}-\frac{\beta L}{2}.
\]
Insert this inequality into the coefficients of $\|x_{k+1}-x_k\|^{2}$ in (7):
\[
-\frac{1}{2\eta}\|x_{k+1}-x_k\|^{2}+\frac{L}{2}\|x_{k+1}-x_k\|^{2}
\le -\frac{\beta L}{2}\|x_{k+1}-x_k\|^{2}. \tag{8}
\]

\subsection*{Step 8 – Simplify the recursion}
Combining (7) and (8) yields
\[
\Phi(x_{k+1})-\Phi^{\star}
\le
\frac{1}{2\eta}\bigl(\|x^{\star}-x_k\|^{2}
-\|x^{\star}-x_{k+1}\|^{2}\bigr)
-\frac{\beta L}{2}\Bigl(\|x_k-x_{k-1}\|^{2}
+\|x_{k+1}-x_k\|^{2}
+\|x_{k+1}-x^{\star}\|^{2}\Bigr). \tag{9}
\]

Discarding the non‑positive last term gives a clean inequality
\[
\Phi(x_{k+1})-\Phi^{\star}
\le
\frac{1}{2\eta}\bigl(\|x^{\star}-x_k\|^{2}
-\|x^{\star}-x_{k+1}\|^{2}\bigr). \tag{10}
\]

\subsection*{Step 9 – Telescoping sum}
Sum \eqref{10} over $k=0,\dots,K-1$:
\[
\sum_{k=0}^{K-1}\bigl(\Phi(x_{k+1})-\Phi^{\star}\bigr)
\le
\frac{1}{2\eta}\bigl(\|x^{\star}-x_0\|^{2}
-\|x^{\star}-x_{K}\|^{2}\bigr)
\le\frac{\|x_0-x^{\star}\|^{2}}{2\eta}. \tag{11}
\]

\subsection*{Step 10 – From average to guarantee}
Since each term on the left of (11) is non‑negative,
\[
K\bigl(\Phi(\bar{x}_{K})-\Phi^{\star}\bigr)
\le
\sum_{k=0}^{K-1}\bigl(\Phi(x_{k+1})-\Phi^{\star}\bigr)
\le\frac{\|x_0-x^{\star}\|^{2}}{2\eta},
\]
which is exactly the bound claimed in Theorem \ref{thm:main}. \hfill $\blacksquare$

---------------------------------------------------------------------

\section*{Rate Derivation (With Constants)}
The theorem yields the explicit constant
\[
C:=\frac{\|x_0-x^{\star}\|^{2}}{2\eta},
\qquad
\Phi(\bar{x}_{K})-\Phi^{\star}\le \frac{C}{K}.
\]
If $f$ is $L$‑smooth and we choose the largest admissible stepsize $\eta=\frac{1}{L(1+\beta)}$, then
\[
C = \frac{L(1+\beta)}{2}\,\|x_0-x^{\star}\|^{2}.
\]

---------------------------------------------------------------------

\section*{Special Cases}
\begin{enumerate}
\item \textbf{No momentum ($\beta=0$).} The update reduces to the classical proximal gradient method and the bound coincides with the standard $O(1/K)$ rate.
\item \textbf{Purely smooth case ($g\equiv0$).} The proximal step becomes the identity, and the algorithm coincides with gradient descent equipped with a momentum‑augmented gradient. The same $O(1/K)$ bound holds.
\item \textbf{Strongly convex $f+g$ (parameter $\mu>0$).} Adding the strong convexity inequality 
$f(y)\ge f(x)+\langle\nabla f(x),y-x\rangle+\frac{\mu}{2}\|y-x\|^{2}$
to the analysis yields a linear convergence rate:
\[
\Phi(x_k)-\Phi^{\star}\le\bigl(1-\eta\mu\bigr)^{k}\bigl(\Phi(x_0)-\Phi^{\star}\bigr).
\]
\end{enumerate}

---------------------------------------------------------------------

\begin{thebibliography}{9}
\bibitem{parikh2014proximal}
N. Parikh and S. Boyd, ``Proximal Algorithms,'' \emph{Foundations and Trends in Optimization}, vol. 1, no. 3, pp. 127–239, 2014.
\end{thebibliography}

\end{document}